% GENERATED FILE. Edit canonical lessons, then run npm run generate:books.
% canonical-source-hash: fnv1a64:1528dccc70d56e45
% canonical-lessons: TE-C323-kante, TE-C323-kalam, TE-C323-samvatsaraluga, TE-R323-comparing-and-time

\chapter{Bigger, Best and How Long}
\label{ch:te-bigger-best-and-how-long}
\hlchaptermodality{323}

\begin{chapteropening}
\textbf{By the end of this chapter:} \emph{I can compare two things with kaṇṭē, pick the best with anniṭikaṇṭē, and say how long with a span and -gā, nuñci and varaku.}

Complete TE-R323-comparing-and-time: compare, pick the best, and say how long something has gone on and since when.
\end{chapteropening}

\section[\texorpdfstring{\te{కంటే} (kaṇṭē)}{కంటే (kaṇṭē)}]{\te{కంటే} (kaṇṭē) — bigger than, best of all}

\label{lesson:TE-C323-kante}



\begin{quote}
Say \textquotedblleft{}big\textquotedblright{} (\emph{pedda}), \textquotedblleft{}more\textquotedblright{} (\emph{ekkuva}) and \textquotedblleft{}I like Telugu\textquotedblright{} (\emph{nāku telugu iṣṭaṁ}). Then \textquotedblleft{}I would like to go to my home town\textquotedblright{}. (\emph{nāku mā ūriki veḷḷālani undi}.)
\end{quote}

\subsection*{You'll want to know: X \te{కంటే}}
\begin{quote}
\textbf{\te{హైదరాబాదు} \te{విజయవాడ} \te{కంటే} \te{పెద్దది}.}
\end{quote}

\begin{quote}
\emph{haidarābādu vijayavāḍa kaṇṭē peddadi.} — \textquotedblleft{}Hyderabad is bigger than Vijayawada.\textquotedblright{}
\end{quote}

\begin{quote}
\textbf{\te{నాకు} \te{టీ} \te{కంటే} \te{కాఫీ} \te{ఎక్కువ} \te{ఇష్టం}.}
\end{quote}

\begin{quote}
\emph{nāku ṭī kaṇṭē kāphī ekkuva iṣṭaṁ.} — \textquotedblleft{}I like coffee more than tea.\textquotedblright{}
\end{quote}

The thing you measure against comes with \textbf{\te{కంటే}}, \textquotedblleft{}than\textquotedblright{}, after it; then the quality. The quality does not change: no \textquotedblleft{}bigger\textquotedblright{}, only \textbf{\te{పెద్దది}} (\emph{peddadi}, \textquotedblleft{}a big one\textquotedblright{}). \textbf{\te{ఎక్కువ}}, \textquotedblleft{}more\textquotedblright{}, can stand before it to make the \textquotedblleft{}more\textquotedblright{} plain.

\begin{quote}
\textbf{\te{రైలు} \te{బస్సు} \te{కంటే} \te{వేగంగా} \te{వెళ్తుంది}.}
\end{quote}

\begin{quote}
\emph{railu bassu kaṇṭē vēgaṅgā veḷtundi.} — \textquotedblleft{}The train goes faster than the bus.\textquotedblright{}
\end{quote}

\subsection*{You'll want to know: \te{అన్నిటికంటే}, \te{అందరికంటే}}
\begin{quote}
\textbf{\te{మా} \te{ఊరిలో} \te{ఈ} \te{హోటలు} \te{అన్నిటికంటే} \te{మంచిది}.}
\end{quote}

\begin{quote}
\emph{mā ūrilō ī hōṭalu anniṭikaṇṭē mañcidi.} — \textquotedblleft{}In our town this restaurant is the best.\textquotedblright{}
\end{quote}

\begin{quote}
\textbf{\te{మా} \te{ఇంట్లో} \te{అమ్మమ్మ} \te{అందరికంటే} \te{పెద్దది}.}
\end{quote}

\begin{quote}
\emph{mā iṇṭlō ammamma andarikaṇṭē peddadi.} — \textquotedblleft{}At home, Grandmother is the eldest.\textquotedblright{}
\end{quote}

\textbf{\te{అన్నిటికంటే}} is \textquotedblleft{}than all (things)\textquotedblright{}; \textbf{\te{అందరికంటే}} is \textquotedblleft{}than everyone\textquotedblright{}, for people. Better than all is the best; bigger than everyone is the eldest.

\subsection*{Your turn}
Say these aloud:

\begin{itemize}
\raggedright
  \item \textquotedblleft{}Hyderabad is bigger than Vijayawada\textquotedblright{} — \emph{haidarābādu vijayavāḍa kaṇṭē peddadi}
  \item \textquotedblleft{}I like coffee more than tea\textquotedblright{} — \emph{nāku ṭī kaṇṭē kāphī ekkuva iṣṭaṁ}
  \item \textquotedblleft{}The train goes faster than the bus\textquotedblright{} — \emph{railu bassu kaṇṭē vēgaṅgā veḷtundi}
  \item \textquotedblleft{}In our town this restaurant is the best\textquotedblright{} — \emph{mā ūrilō ī hōṭalu anniṭikaṇṭē mañcidi}
  \item two people you know, compared: who is older, who is the eldest
\end{itemize}

\begin{tcolorbox}[breakable,colback=teal!4,colframe=teal!35!black,title={Before you move on}]
\textquotedblleft{}Than\textquotedblright{}? (\textbf{\te{కంటే}}.) Where does it stand? (\textbf{After the thing you measure against.}) \textquotedblleft{}Best of all\textquotedblright{}? (\textbf{\te{అన్నిటికంటే} \te{మంచిది}}.) \textquotedblleft{}Than everyone\textquotedblright{}? (\textbf{\te{అందరికంటే}}.)
\end{tcolorbox}
\section[\texorpdfstring{\te{కాలం} (kālaṁ)}{కాలం (kālaṁ)}]{\te{కాలం} (kālaṁ) — a span of time}

\label{lesson:TE-C323-kalam}



\begin{quote}
Before the new one: say \textquotedblleft{}time\textquotedblright{} (\emph{samayaṁ}). Then the four seasons you learned long ago. (\emph{vasanta ṛtuvu}, \emph{vēsavi}, \emph{vānākālaṁ}, \emph{śītākālaṁ}.)
\end{quote}

\subsection*{You'll want to know: \te{కాలం}}
\textbf{\te{కాలం}} — \emph{kālaṁ} — \textquotedblleft{}a span of time; a period, an age\textquotedblright{}.

\begin{quote}
\textbf{\te{చాలా} \te{కాలం} \te{తరువాత} \te{నేను} \te{మా} \te{ఊరికి} \te{వెళ్ళాను}.}
\end{quote}

\begin{quote}
\emph{cālā kālaṁ taruvāta nēnu mā ūriki veḷḷānu.} — \textquotedblleft{}After a long time I went to my home town.\textquotedblright{}
\end{quote}

\begin{quote}
\textbf{\te{మీరు} \te{ఇక్కడ} \te{ఎంత} \te{కాలం} \te{ఉంటారు}?}
\end{quote}

\begin{quote}
\emph{mīru ikkaḍa enta kālaṁ uṇṭāru?} — \textquotedblleft{}How long will you stay here?\textquotedblright{}
\end{quote}

\textbf{\te{సమయం}} is time on the clock: is there time, what time is it. \textbf{\te{కాలం}} is time as a stretch: how long, a long while, a season. You have said it already without noticing it: \textbf{\te{వానాకాలం}} is \textquotedblleft{}the rains' time\textquotedblright{}, the rainy season, and \textbf{\te{శీతాకాలం}} is \textquotedblleft{}the cold time\textquotedblright{}, winter.

\begin{cousinweb}
\textbf{\te{కాలం}} (\emph{kālaṁ}) is borrowed from Sanskrit \emph{kāla}, \textquotedblleft{}time\textquotedblright{} (CDIAL 3084), with Telugu's noun ending \textbf{-\te{ం}}. The same \emph{kāla} is the first half of \textbf{\te{కాలపట్టిక}}, \textquotedblleft{}a timetable\textquotedblright{}.
\end{cousinweb}

\subsection*{Your turn}
Say these aloud:

\begin{itemize}
\raggedright
  \item \emph{kālaṁ}
  \item \textquotedblleft{}After a long time I went to my home town\textquotedblright{} — \emph{cālā kālaṁ taruvāta nēnu mā ūriki veḷḷānu}
  \item \textquotedblleft{}How long will you stay here?\textquotedblright{} — \emph{mīru ikkaḍa enta kālaṁ uṇṭāru?}
  \item the two seasons that end in \emph{kālaṁ}
\end{itemize}

\begin{tcolorbox}[breakable,colback=teal!4,colframe=teal!35!black,title={Before you move on}]
What does \textbf{\te{కాలం}} mean? (\textbf{A span of time.}) And \textbf{\te{సమయం}}? (\textbf{Time on the clock.}) \textquotedblleft{}How long?\textquotedblright{} (\textbf{\te{ఎంత} \te{కాలం}?})
\end{tcolorbox}
\section[\texorpdfstring{\te{మూడు} \te{సంవత్సరాలుగా} (mūḍu saṁvatsarālugā)}{మూడు సంవత్సరాలుగా (mūḍu saṁvatsarālugā)}]{\te{మూడు} \te{సంవత్సరాలుగా} (mūḍu saṁvatsarālugā) — how long}

\label{lesson:TE-C323-samvatsaraluga}



\begin{quote}
Say \textquotedblleft{}a year\textquotedblright{} (\emph{saṁvatsaraṁ}), \textquotedblleft{}from\textquotedblright{} (\emph{nuñci}) and \textquotedblleft{}until\textquotedblright{} (\emph{varaku}). Then \textquotedblleft{}how long?\textquotedblright{} (\emph{enta kālaṁ?})
\end{quote}

\subsection*{You'll want to know: a span + -\te{గా}, 'for … now'}
\begin{quote}
\textbf{\te{మూడు} \te{సంవత్సరాలుగా} \te{నేను} \te{తెలుగు} \te{నేర్చుకుంటున్నాను}.}
\end{quote}

\begin{quote}
\emph{mūḍu saṁvatsarālugā nēnu telugu nērcukuṇṭunnānu.} — \textquotedblleft{}I have been learning Telugu for three years.\textquotedblright{}
\end{quote}

\begin{quote}
\textbf{\te{రెండు} \te{రోజులుగా} \te{వర్షం} \te{పడుతోంది}.}
\end{quote}

\begin{quote}
\emph{reṇḍu rōjulugā varṣaṁ paḍutōndi.} — \textquotedblleft{}It has been raining for two days.\textquotedblright{}
\end{quote}

The span takes \textbf{-\te{గా}}, and the verb stays in the \textbf{-\te{తున్న}-} form: it began then and is still going on. English says \textquotedblleft{}have been\textquotedblright{}; Telugu says \textquotedblleft{}am -ing\textquotedblright{}. The plural of \textbf{\te{సంవత్సరం}} is \textbf{\te{సంవత్సరాలు}}; of \textbf{\te{రోజు}}, \textbf{\te{రోజులు}}.

\begin{quote}
\textbf{\te{ఎంత} \te{కాలంగా} \te{మీరు} \te{ఇక్కడ} \te{ఉంటున్నారు}?}
\end{quote}

\begin{quote}
\emph{enta kālaṅgā mīru ikkaḍa uṇṭunnāru?} — \textquotedblleft{}How long have you been living here?\textquotedblright{}
\end{quote}

A span that is over takes no \textbf{-\te{గా}}:

\begin{quote}
\textbf{\te{నేను} \te{మూడు} \te{గంటలు} \te{నడిచాను}.}
\end{quote}

\begin{quote}
\emph{nēnu mūḍu gaṇṭalu naḍicānu.} — \textquotedblleft{}I walked for three hours.\textquotedblright{}
\end{quote}

\subsection*{You'll want to know: \te{నుంచి} … \te{వరకు}}
\begin{quote}
\textbf{\te{నేను} \te{ఉదయం} \te{తొమ్మిది} \te{నుంచి} \te{సాయంత్రం} \te{ఐదు} \te{వరకు} \te{పని} \te{చేస్తాను}.}
\end{quote}

\begin{quote}
\emph{nēnu udayaṁ tommidi nuñci sāyantraṁ aidu varaku pani cēstānu.} — \textquotedblleft{}I work from nine in the morning until five in the evening.\textquotedblright{}
\end{quote}

\begin{quote}
\textbf{\te{నిన్నటి} \te{నుంచి} \te{నాకు} \te{జ్వరం} \te{ఉంది}.}
\end{quote}

\begin{quote}
\emph{ninnaṭi nuñci nāku jvaraṁ undi.} — \textquotedblleft{}I have had a fever since yesterday.\textquotedblright{}
\end{quote}

\textbf{\te{నుంచి}} marks where a span starts, \textbf{\te{వరకు}} where it stops. \textbf{\te{నిన్నటి}} (\emph{ninnaṭi}, \textquotedblleft{}yesterday's\textquotedblright{}) is the form \textbf{\te{నిన్న}} takes before an ending: \textquotedblleft{}from yesterday\textquotedblright{}, so \textquotedblleft{}since yesterday\textquotedblright{}.

\subsection*{Your turn}
Say these aloud:

\begin{itemize}
\raggedright
  \item \textquotedblleft{}I have been learning Telugu for three years\textquotedblright{} — \emph{mūḍu saṁvatsarālugā nēnu telugu nērcukuṇṭunnānu}
  \item \textquotedblleft{}I walked for three hours\textquotedblright{} — \emph{nēnu mūḍu gaṇṭalu naḍicānu}
  \item \textquotedblleft{}I work from nine in the morning until five in the evening\textquotedblright{} — \emph{nēnu udayaṁ tommidi nuñci sāyantraṁ aidu varaku pani cēstānu}
  \item how long you have been learning Telugu, and since when
\end{itemize}

\begin{tcolorbox}[breakable,colback=teal!4,colframe=teal!35!black,title={Before you move on}]
\textquotedblleft{}For two years now\textquotedblright{}? (\textbf{\te{రెండు} \te{సంవత్సరాలుగా}}.) Which form does the verb keep? (\textbf{-\te{తున్న}-}, still going on.) \textquotedblleft{}Since yesterday\textquotedblright{}? (\textbf{\te{నిన్నటి} \te{నుంచి}}.) \textquotedblleft{}Until five\textquotedblright{}? (\textbf{\te{ఐదు} \te{వరకు}}.)
\end{tcolorbox}
\section[(review)]{Review — comparing and time}

\label{lesson:TE-R323-comparing-and-time}



\begin{quote}
Say the five words this core took its letters from. (\emph{īlōgā}, \emph{ōṭu}, \emph{ūru}, \emph{auṣadhaṁ}, \emph{duḥkhaṁ}.) Then say \textquotedblleft{}a span of time\textquotedblright{}. (\emph{kālaṁ}.)

\begin{itemize}
\raggedright
  \item \emph{Write it:} the five letters, \textbf{\te{ఈ}}, \textbf{\te{ఓ}}, \textbf{\te{ఊ}}, \textbf{\te{ఔ}}, \textbf{\te{ః}}, once each from memory
\end{itemize}
\end{quote}

\subsection*{Your turn}
\begin{itemize}
\raggedright
  \item \emph{Say it:} My home town is smaller than Hyderabad, but for me it is the best of all.
  \item \emph{Say it:} I have been living here for two years.
  \item \emph{Say it:} After a long time I went to the village. I felt very happy.
  \item \emph{Say it:} I work from morning until evening. I would like to go home.
  \item \emph{Write it:} \textbf{\te{ఈలోగా}}, \textbf{\te{ఓటు}}, \textbf{\te{ఊరు}}, \textbf{\te{ఔషధం}}, \textbf{\te{దుఃఖం}}
\end{itemize}

\begin{tcolorbox}[breakable,colback=teal!4,colframe=teal!35!black,title={Before you move on}]
\textquotedblleft{}Than\textquotedblright{}? (\textbf{\te{కంటే}}.) \textquotedblleft{}The best\textquotedblright{}? (\textbf{\te{అన్నిటికంటే} \te{మంచిది}}.) \textquotedblleft{}For three years now\textquotedblright{}? (\textbf{\te{మూడు} \te{సంవత్సరాలుగా}}.) \textquotedblleft{}From … until\textquotedblright{}? (\textbf{\te{నుంచి} … \te{వరకు}}.) \textquotedblleft{}How long?\textquotedblright{} (\textbf{\te{ఎంత} \te{కాలం}?})
\end{tcolorbox}
